\section{The oriented twist and a finite matrix comparison}
\label{sec:twist-matrix-comparison}

We determine the conjugating proword directly from two pentagon
identities, then compute the complete two-dimensional homology
matrix of the twist. The matrix gives the same degree--163 scalar
test as before and explains a limitation for central commutators.

Use the conventions
\[
 \operatorname{Inn}(g)(h)=ghg^{-1},\qquad
 x_{ij}=\sigma_{j-1}\cdots\sigma_{i+1}\sigma_i^2
                \sigma_{i+1}^{-1}\cdots\sigma_{j-1}^{-1}
                    \quad(i<j).
\]
Work in the pure five-point sphere mapping class group and its
profinite completion, and abbreviate
\[
 u=x_{12},\quad v=x_{23},\quad
 x=x_{14},\quad y=x_{24},\quad z=x_{34},\quad
 w=x_{45},\quad t=x_{51}.
\]
Then \(xyzw=1\). Set
\[
 c=x'_{35}=\sigma_4^{-1}\sigma_3^2\sigma_4.
\]
The two sphere identities used below are
\begin{equation}
 c=uxy,\qquad t=vyz.                                  \label{eq:twist-sphere-complements}
\end{equation}
They identify the boundary twist of three indicated punctures with
the twist of the complementary two punctures. In these conventions
they are equation~(18) and a cyclic instance of equation~(16) in
\cite{LS2015Twists}; the braid-word definition is its equation~(14).
These are relations in the sphere mapping class group and do not
use any GT action formula.

\subsection{Pure-braid conjugation on the fiber}

\begin{lemma}\label{lem:twist-direct-pure-braid}
On the fiber \(\widehat F_3=\langle x,y,z\rangle\),
\begin{align*}
 \operatorname{Inn}(u)(x)&=(xy)^{-1}x(xy)=y^{-1}xy,\\
 \operatorname{Inn}(u)(y)&=(xy)^{-1}y(xy)=y^{-1}x^{-1}yxy,\\
 \operatorname{Inn}(u)(z)&=z,
\end{align*}
and
\begin{align*}
 \operatorname{Inn}(v)(x)&=x,\\
 \operatorname{Inn}(v)(y)&=(yz)^{-1}y(yz)=z^{-1}yz,\\
 \operatorname{Inn}(v)(z)&=(yz)^{-1}z(yz)=z^{-1}y^{-1}zyz.
\end{align*}
In particular both base generators fix \(xyz\), and commute
with \(w=(xyz)^{-1}\).
\end{lemma}

\begin{proof}
In the three-strand braid group put \(a=\sigma_1\),
\(b=\sigma_2\). The pure elements are
\(a^2,ba^2b^{-1},b^2\). Their product is
\[
 a^2(ba^2b^{-1})b^2=a^2ba^2b
      =a(aba)ab=a(bab)ab=(ab)^3,
\]
where the middle equality is the braid relation \(aba=bab\).
This is central: if \(\Delta=aba=bab\), then
\(\Delta a=b\Delta\) and \(\Delta b=a\Delta\), and
\((ab)^3=\Delta^2\). Hence conjugation by \(a^2\) on the
other two pure generators is conjugation by the inverse of their
product.

For the strands \(1,2,4\), conjugate this three-strand calculation
by \(\sigma_3\). The base \(\sigma_1^2=u\) is fixed,
because \(\sigma_1\) and \(\sigma_3\) commute. The other two
pure elements become \(x\) and \(y\). Thus \(uxy\) commutes
with \(u,x,y\), proving the first two formulas. The last follows
from \([\sigma_1,\sigma_3]=1\).

The same three-strand argument with \(a=\sigma_2,b=\sigma_3\)
shows that \(vyz\) is central in \(\langle v,y,z\rangle\),
proving the last two formulas for \(v\). To verify \([v,x]=1\)
directly, put \(b=\sigma_2\), \(d=\sigma_3\), and
\(q=db\). The braid relation gives
\[
 d^{-1}bd=bdb^{-1},\qquad q^{-1}bq=d.
\]
Since \(x=q\sigma_1^2q^{-1}\) and \(d\) commutes with
\(\sigma_1\), it follows that \(b\), and hence \(v=b^2\),
commutes with \(x\). These identities hold already in the
discrete braid groups, so they persist in the indicated quotient
and completion. The product assertion follows immediately.
\end{proof}

\subsection{Two pentagons determine the orientation}

Let \((1,f)\) be a full GT element. Use the peripheral action
based at infinity and the four block-substitution actions
\(F^{1,2},F^{2,3},F^{12,3},F^{1,23}\) of equations
(9)--(12) of \cite{LS2015Twists}, with the disclosed correction of the first
argument in \(F^{2,3}(z)\) from \(y\) to \(z\).
Explicitly, with \(a=f(x,(xy)^{-1})\),
\(b=f(y,(xy)^{-1})\), the first action sends
\((x,y,z)\) to \((axa^{-1},byb^{-1},z)\). The second
is the same two-generator construction on \((y,z)\), fixing
\(x\). The third conjugates \(x,y\) by \(f(xy,w)\)
and \(z\) by \(f(z,w)\); the fourth conjugates \(x\)
by \(f(x,w)\) and \(y,z\) by \(f(yz,w)\).
They preserve \(xyz\); moreover \(F^{1,2}\) fixes \(xy,z\),
and \(F^{2,3}\) fixes \(x,yz\). These are the elementary
product identities of Lemma~3.1 of \cite{LS2015Twists}, derived from the two- and
three-cycle relations, rather than from the twist theorem.

Set
\[
 L=F^{1,2}F^{12,3},\qquad R=F^{2,3}F^{1,23},
 \qquad d=f(v,u)=f(u,v)^{-1}.
\]
Products of automorphisms mean composition, with the rightmost
factor applied first.

\begin{proposition}\label{prop:twist-oriented-pentagon}
With these conventions,
\begin{equation}
 L=\operatorname{Inn}(f(v,u))R,
 \qquad
 \operatorname{Inn}(f(x_{23},x_{12}))
 =F^{1,2}F^{12,3}(F^{1,23})^{-1}(F^{2,3})^{-1}.
                                                               \label{eq:twist-oriented}
\end{equation}
The equalities are automorphisms of the rank-three fiber.
\end{proposition}

\begin{proof}
We will use the elementary central-factor rule: if \(a,b,e\)
are group elements with \(e\) commuting with \(a,b\), then
\[
 f(ae,b)=f(a,b)=f(a,be).
\]
For discrete commutator words this follows by collecting the powers
of the central factor; their exponent sum is zero. Continuity gives
the rule for a proword in the closed commutator subgroup.

The block definitions and their product preservation give
\[
 L(x)=\operatorname{Inn}\bigl(f(xy,w)f(x,(xy)^{-1})\bigr)(x),
 \qquad L(z)=\operatorname{Inn}(f(z,w))(z),
\]
and
\[
 R(x)=\operatorname{Inn}(f(x,w))(x),\qquad
 R(z)=\operatorname{Inn}\bigl(f(yz,w)f(z,(yz)^{-1})\bigr)(z).
\]
Because \(u\) commutes with \(xy,w\), and \(c=uxy\)
commutes with \(u,x\), the first conjugator for \(L(x)\) is
\[
 a=f(c,w)f(x,u).
\]
Likewise \(v\) commutes with \(yz,w\), while \(t=vyz\)
commutes with \(v,z\), so the conjugator for \(R(z)\) is
\[
 a_z=f(t,w)f(z,v).
\]

The defining pentagon is
\begin{equation}
 f(u,v)f(z,w)f(t,u)f(v,z)f(w,t)=1.                    \label{eq:twist-pentagon}
\end{equation}
Conjugate this identity by \(\sigma_4^{-1}\). The braid-word
definitions show that \(u,v,w\) are fixed, \(z\) goes to
\(c=\sigma_4^{-1}\sigma_3^2\sigma_4\), and \(t\) goes
to \(x\), since \(t=x_{15}=\sigma_4x_{14}\sigma_4^{-1}\).
Thus
\begin{equation}
 f(u,v)f(c,w)f(x,u)f(v,c)f(w,x)=1.                    \label{eq:twist-conjugated-pentagon}
\end{equation}
Put \(b=f(x,w)\). By duality, \eqref{eq:twist-conjugated-pentagon} says
\[
 a=d\,b\,f(c,v).
\]
Both \(c\) and \(v\) commute with \(x\), so the final
factor commutes with \(x\). Hence
\[
 L(x)=\operatorname{Inn}(d)R(x).
\]

Next put \(e=f(z,w)\) and \(q=f(t,u)\). Since
\(a_z^{-1}=f(v,z)f(w,t)\), \eqref{eq:twist-pentagon} gives
\[
 e\,q=d\,a_z.
\]
Both \(t\) and \(u\) commute with \(z\), so \(q\)
commutes with \(z\). It follows that
\[
 L(z)=\operatorname{Inn}(d)R(z).
\]
Finally \(L\), \(R\), and \(\operatorname{Inn}(d)\)
all fix \(xyz\). Equality on \(x,z,xyz\) therefore gives
equality on \(y\) as well. This proves \eqref{eq:twist-oriented}.
\end{proof}

\paragraph{Comparison of conventions.}
With \(\operatorname{Inn}(g)(h)=ghg^{-1}\), the proposition
has \(f(x_{23},x_{12})\) on the left. This reverses the two
arguments in the printed equation~(8) of \cite{LS2015Twists}.
The proof derives the orientation directly from the defining
pentagon and the explicit block actions. It does not use the
standard-action display~(26) or Propositions~2.2--2.3 of that
source.

\subsection{The exact finite matrix identity}

Let \(O=\mathbf Z_2[\zeta_{163}]\), choose a primitive
\(163\)-rd root \(t\) in \(O/4O\), and put \(q=1+t\).
For units \(u\), define
\[
 C_t(u)=\begin{pmatrix}u&-(u-1)/q\\0&1\end{pmatrix},\qquad
 D_t(u)=\begin{pmatrix}1&0\\-t(u-1)/q&u\end{pmatrix}.
\]
For the local system sending each of \(x,y,z\) to \(t\),
take the homology basis
\[
 v_1=(1,-1,0),\qquad v_2=(0,1,-1).
\]
The direct conjugation lemma gives
\begin{equation}
 A_t=C_t(t^{-2}),\qquad B_t=D_t(t^{-2})                \label{eq:twist-base-matrices}
\end{equation}
for \(\operatorname{Inn}(u),\operatorname{Inn}(v)\).
Indeed the standard Artin half-twists have matrices
\[
 S_1=\begin{pmatrix}-t&1\\0&1\end{pmatrix},\qquad
 S_2=\begin{pmatrix}1&0\\t&-t\end{pmatrix},
\]
and the lemma identifies the pure-braid actions with their inverse
squares. The equalities \(S_1^2=C_t(t^2)\),
\(S_2^2=D_t(t^2)\), and the multiplicativity of \(C_t,D_t\)
give \eqref{eq:twist-base-matrices}.

Write
\[
 a=B_f(t,t),\qquad b=B_f(t^2,t)=B_f(t,t^2),\qquad k=a/b.
\]
The four block matrices, in the same column convention, are
\[
 M_{1,2}=C_t(a),\quad M_{2,3}=D_t(a),\quad
 M_{12,3}=bC_t(b^{-1}),\quad
 M_{1,23}=bD_t(b^{-1}).
\]
For completeness, a two-generator peripheral action fixes the product
chain \((1,t)\) and acts by its beta scalar on the cycle
\((1,-1)\). These form a basis since \(1+t\) is a unit.
This yields the first two matrices. The third action fixes \(v_1\)
and acts by \(b\) on \(v_1+(1+t)v_2\), arising from the
subgroup generated by \(xy,z\). The fourth fixes \(v_2\)
and acts by \(b\) on \((1+t)v_1+t v_2\), arising from
\(x,yz\). These properties give the last two matrices.

Proposition~\ref{prop:twist-oriented-pentagon} therefore gives the fully normalized identity
\begin{equation}
 f(B_t,A_t)=C_t(k)D_t(k)^{-1},\qquad
 f(A_t,B_t)=D_t(k)C_t(k)^{-1}.                        \label{eq:twist-finite-word}
\end{equation}
The second equality is duality. Every proword evaluation here is
in one fixed finite matrix group.

The first matrix on the right is
\[
 \begin{pmatrix}
 k-\dfrac{t(k-1)^2}{q^2k}&-\dfrac{k-1}{qk}\\[4pt]
 \dfrac{t(k-1)}{qk}&k^{-1}
 \end{pmatrix}.
\]
If \(k=1+2u\pmod{4O}\), it equals
\begin{equation}
 I+2u\begin{pmatrix}1&-q^{-1}\\tq^{-1}&-1\end{pmatrix}.
                                                               \label{eq:twist-first-layer}
\end{equation}
Its inverse has the same value modulo four. In particular its
upper-left entry is \(1+2u\), so this coefficient is not lost.

For \(r\in R_{\rm global}\), let \(n=2^{81}+1\) and use
the proword of the GT group power \(r^n\), not the ordinary word
power of \(f_r\). The beta ratio of \(r^n\) is \(k_r^n\)
and is principal. The gamma-ratio identity of Proposition~\ref{prop:163-mixed-beta-residual} gives
\[
 \Psi(k_r^n)=\Psi(h_{163}(r)).
\]
Consequently \eqref{eq:twist-finite-word}--\eqref{eq:twist-first-layer}
are an exact finite matrix representation of the same unresolved
degree--163 scalar coordinate.

\begin{corollary}\label{cor:twist-finite-scalar-and-commutator}
For an acyclotomic full GT element, the finite matrices satisfy
\eqref{eq:twist-finite-word}. For \(r\in R_{\rm global}\),
the first layer after the GT group power \(2^{81}+1\) determines
\(\Psi(h_{163}(r))\). For any two multiplier-one elements
\(r_1,r_2\), one has
\[
 f_{[r_1,r_2]}(A_t,B_t)=1.
\]
Thus this specialization cannot detect the commutator value of the
relative quaternion obstruction.
\end{corollary}
\begin{proof}
The matrix calculation and gamma comparison just proved give the
first two assertions. On the multiplier-one subgroup, each beta evaluation, and hence the
ratio \(k_r\), is multiplicative in the GT group law. Thus
\(k_{[r_1,r_2]}=1\), and \eqref{eq:twist-finite-word} gives
\[
 f_{[r_1,r_2]}(A_t,B_t)=1
\]
for multiplier-one elements \(r_1,r_2\). This proves the last assertion.
\end{proof}

The matrix identity remains an equivalent test for the full scalar
image if an additional full-GT identity controls these proword
values. Such an identity is not proved here.
